\documentclass[conference]{IEEEtran}
\IEEEoverridecommandlockouts

\usepackage{cite}
\usepackage{amsmath,amssymb}
\usepackage{graphicx}
\usepackage{xcolor}
\usepackage{url}

% ---------------------------------------------------------------------------
% Helper macros
% \todo{...}  -> red text for things the authors must fill in
% \figph[width]{description}{file} -> inserts figures/<file>.png when it exists
%               (generate them with: python src/make_figures.py); otherwise it
%               shows a framed placeholder so the document always compiles.
% ---------------------------------------------------------------------------
\newcommand{\todo}[1]{\textcolor{red}{\textbf{[TODO: #1]}}}
\newcommand{\figph}[3][\columnwidth]{%
  \IfFileExists{figures/#3.png}{%
    \includegraphics[width=#1]{figures/#3.png}}{%
    \fbox{\parbox[c][2.6cm][c]{\dimexpr#1-2\fboxsep-2\fboxrule\relax}{\centering\small
    \textbf{Insert figure:} #2\\[2pt] \texttt{figures/#3.png}}}}}

\begin{document}

\title{Data Preparation and Supervised Baselines for Failure Prediction in CartPole Trajectories}

\author{\IEEEauthorblockN{Miguel Angel Torres Ramirez}
\IEEEauthorblockA{Dept. of Computer Engineering\\
Universidad Distrital Francisco Jos\'e de Caldas\\
Email: migatorresr@udistrital.edu.co}
\and
\IEEEauthorblockN{Cristian David Romero Gil}
\IEEEauthorblockA{Dept. of Computer Engineering\\
Universidad Distrital Francisco Jos\'e de Caldas\\
Email: cdromerog@udistrital.edu.co}
\and
\IEEEauthorblockN{Omar Esteban Peña Sanchez}
\IEEEauthorblockA{Dept. of Computer Engineering\\
Universidad Distrital Francisco Jos\'e de Caldas\\
Email: oepenas@udistrital.edu.co}}

\maketitle

\begin{abstract}
This report presents the data preparation pipeline and the supervised baselines
for a binary failure-prediction task built on CartPole trajectories
(35,874 samples from 400 episodes). The label marks the last five steps of every
episode that ends in failure, which yields a 17.2:1 class imbalance. We describe the
exploratory analysis, an episode-level split that prevents leakage between
consecutive samples, a physics-informed feature set, and three traditional
classifiers: Logistic Regression, Random Forest and Gradient Boosting. On the held-out
test episodes, Gradient Boosting obtains the best results (F1 = 0.860, PR-AUC = 0.935),
while Logistic Regression only becomes competitive (ROC-AUC 0.981 versus 0.542) after
feature engineering. These results establish the reference performance for the
deep learning models of the next workshops.
\end{abstract}

\begin{IEEEkeywords}
CartPole, failure prediction, class imbalance, exploratory data analysis, Random Forest, Gradient Boosting, baselines.
\end{IEEEkeywords}

%% ======================================================================
\section{Introduction}
Anticipating that a control system is about to fail is a natural supervised learning
problem: given the current state of the system, a classifier must raise an alarm a few
steps before the failure occurs. In this workshop we study this problem on trajectories of
the CartPole system \cite{barto1983}, where a pole must be kept upright on a moving cart.

The central question is: \emph{what does the data look like, and what performance can
well-engineered traditional models achieve?} The answer fixes the difficulty of the
problem, the degree of class imbalance and the reference metrics against which the deep
learning models of later workshops will be compared.

The contributions of this report are: (i) a verified and characterised dataset;
(ii) a leakage-free preprocessing pipeline based on an episode-level split;
(iii) three baselines evaluated with metrics suited to imbalanced data; and
(iv) two controlled ablations that isolate the effect of feature engineering and of data
augmentation.

The rest of the paper is organised as follows: Section~II describes the dataset and its
exploratory analysis; Section~III presents the preprocessing pipeline; Section~IV details
the baseline models and the evaluation protocol; Section~V reports and discusses the
results; and Section~VI concludes.

%% ======================================================================
\section{Dataset and Exploratory Data Analysis}

\subsection{Dataset Description}
The dataset contains one row per time step of each simulated episode. Each row stores the
episode identifier, the time index \texttt{t}, the cart position \texttt{x}, cart velocity
\texttt{x\_dot}, pole angle \texttt{theta}, pole angular velocity \texttt{theta\_dot}, the
\texttt{action} taken (0/1) and the binary target \texttt{label}. The task is binary
classification. The dataset was generated by rolling out the CartPole-v1 environment
(OpenAI Gym 0.26.2) with a mixed policy (noisy expert plus random actions), with all
seeds (Python, NumPy, Gym) fixed to 42.
Each state is labelled 1 if the episode physically fails within the next five steps
($|\theta|$ beyond 12$^\circ$ or $|x|>2.4$ m). The environment is free and pip-installable,
with no license restrictions for this use.

Integrity was verified with the SHA256 checksum
(\texttt{ec067d5d\dots f53f8}), which matches the expected value. Table~\ref{tab:dataset}
summarises the dataset, as printed by \texttt{dataset\_info.py}.

\begin{table}[t]
\caption{DATASET SUMMARY}
\label{tab:dataset}
\centering
\begin{tabular}{|l|l|}
\hline
\textbf{Property} & \textbf{Value} \\ \hline
File size & 1.99 MB \\ \hline
Rows / episodes & 35,874 / 400 \\ \hline
Missing values / duplicates & 0 / 0 \\ \hline
Class 0 (safe) & 33,904 (94.51\%) \\ \hline
Class 1 (failure imminent) & 1,970 (5.49\%) \\ \hline
Imbalance ratio & 17.2 : 1 \\ \hline
Episode length (min / median / max) & 8 / 34.5 / 500 \\ \hline
\end{tabular}
\end{table}

\subsection{Target Definition and Class Imbalance}
The label equals 1 exactly in the last five steps of the 394 episodes that end in failure.
The remaining six episodes reach the 500-step limit and contain no positive sample. With a
17.2:1 ratio, a trivial classifier that always predicts ``safe'' already reaches 94.5\%
accuracy, so accuracy cannot be used as a decision metric.

\begin{figure*}[t]
\centering
\figph[0.95\textwidth]{class counts, episode-length histogram and label vs. steps to the end}{fig-target-episodes}
\caption{Target distribution, episode-length distribution (log scale) and label probability as a function of the steps remaining before the end of the episode.}
\label{fig:target}
\end{figure*}

\subsection{Feature Statistics and Data Quality}
The four state variables are well bounded: \texttt{x} $\in[-2.40,2.40]$,
\texttt{theta} $\in[-0.209,0.209]$ rad, and the velocities stay within roughly $\pm 3$.
No missing values, duplicates or corrupted rows were found, so 100\% of the samples are
usable. The IQR rule flags 26.4\% of the \texttt{x} values as outliers, but this reflects
the wide spread of the cart position, not erroneous data; the other state variables show
between 3.0\% and 4.3\% flagged values. All stored states lie inside the standard failure
limits ($|\theta|<0.2095$ rad, $|x|<2.4$ m); the terminal state is not recorded.

\subsection{Feature Relationships}
The Pearson correlation between each raw feature and the label is very weak
(between $-0.036$ and $0.015$), which indicates that no raw variable separates the classes
linearly. Fig.~\ref{fig:dist} shows the class-conditional distributions and Fig.~\ref{fig:corr}
the correlation matrix.

\begin{figure*}[t]
\centering
\figph[0.95\textwidth]{class-conditional histograms of the four state variables}{fig-distributions}
\caption{Distribution of the state variables by class (density, label 0 = safe, 1 = danger).}
\label{fig:dist}
\end{figure*}

\begin{figure}[t]
\centering
\figph{correlation matrix of the raw features and the label}{fig-corr}
\caption{Correlation matrix of the raw features and the label.}
\label{fig:corr}
\end{figure}

%% ======================================================================
\section{Preprocessing Pipeline}

\subsection{Episode-Level Split}
Consecutive rows of an episode are nearly identical, so a random split by row would place
almost duplicated states in train and test and inflate the metrics. We therefore split the
400 \emph{episodes} into 70/15/15\%, stratified by episode-length quartile because the
lengths are strongly skewed (Table~\ref{tab:split}). The split is stored in
\texttt{data/splits.json} and is never recomputed.

\begin{table}[t]
\caption{EPISODE-LEVEL SPLIT (SEED 42)}
\label{tab:split}
\centering
\begin{tabular}{|l|c|c|c|}
\hline
\textbf{Partition} & \textbf{Episodes} & \textbf{Rows} & \textbf{Positives} \\ \hline
Train & 279 & 24,782 & 5.57\% \\ \hline
Validation & 61 & 5,762 & 5.21\% \\ \hline
Test & 60 & 5,330 & 5.44\% \\ \hline
\end{tabular}
\end{table}

\subsection{Feature Engineering}
Because an episode fails when $|\theta|>0.2095$ rad or $|x|>2.4$ m, the features describe the
distance to those limits and the speed of approach. All of them use only the current and past
steps of the same episode (no information from the future), and the time index \texttt{t} is
excluded. The final 16 features are:
\begin{itemize}
\item \textbf{Raw (5):} \texttt{x}, \texttt{x\_dot}, \texttt{theta}, \texttt{theta\_dot}, \texttt{action}.
\item \textbf{Distance to limits (4):} $|x|$, $|\theta|$, $x$-margin and $\theta$-margin.
\item \textbf{Short-term dynamics (3):} first differences of $x$, $\theta$ and $\dot\theta$.
\item \textbf{Rolling windows (2):} 3-step mean of $|\theta|$ and $|\dot\theta|$.
\item \textbf{Time to limit (2):} seconds to reach each limit assuming constant velocity.
\end{itemize}
The engineered features correlate far more with the label than the raw ones
(e.g., 0.561 for the rolling mean of $|\dot\theta|$ and 0.472 for $|\theta|$).

\subsection{Normalisation, Missing Data and Augmentation}
Features are standardised with z-scores because their ranges differ widely and the linear
model is scale-sensitive; the mean and standard deviation are computed \emph{only} on the
training set and then applied to validation and test. There are no missing values, so no
imputation is needed. As an optional augmentation, applied to the training set only,
we create one noisy copy of each training episode by adding Gaussian noise equal to 2\% of
the standard deviation of each state variable, which emulates sensor imprecision while
keeping the label.

%% ======================================================================
\section{Baseline Models and Evaluation Protocol}

\subsection{Models}
Three traditional classifiers were trained on the scaled features (scikit-learn
\cite{pedregosa2011}):
\begin{itemize}
\item \textbf{Baseline 1 -- Logistic Regression:} linear model with balanced class weights.
\item \textbf{Baseline 2 -- Random Forest} \cite{breiman2001}: 300 trees, balanced-subsample class weights.
\item \textbf{Baseline 3 -- Gradient Boosting} \cite{friedman2001}: histogram-based implementation, 300 iterations, learning rate 0.1, balanced sample weights.
\end{itemize}

\subsection{Evaluation Protocol}
Models are fitted on the training set. The decision threshold is selected on the validation
set by maximising F1, so the test set is never used for tuning. We report accuracy,
precision, recall, F1, ROC-AUC, PR-AUC and the confusion matrix. For this task a missed
failure (false negative) is the costly error, so \textbf{recall} is the key metric, balanced
by precision to limit false alarms; PR-AUC is preferred over ROC-AUC for summarising
performance under a 17:1 imbalance \cite{saito2015}. Accuracy is reported only for completeness.

%% ======================================================================
\section{Results and Discussion}

\subsection{Baseline Comparison}
Table~\ref{tab:results} reports the test-set results (60 episodes, 5,330 rows, 290 positives)
at the validation threshold. Fig.~\ref{fig:cm} shows the confusion matrices.

\begin{table}[t]
\caption{BASELINE RESULTS ON THE TEST SET (ENGINEERED FEATURES)}
\label{tab:results}
\centering
\begin{tabular}{|l|c|c|c|c|c|}
\hline
\textbf{Model} & \textbf{Prec.} & \textbf{Rec.} & \textbf{F1} & \textbf{ROC-AUC} & \textbf{PR-AUC} \\ \hline
Logistic Regression & 0.669 & 0.703 & 0.686 & 0.981 & 0.791 \\ \hline
Random Forest & 0.812 & 0.852 & 0.832 & 0.994 & 0.923 \\ \hline
\textbf{Gradient Boosting} & 0.834 & \textbf{0.886} & \textbf{0.860} & \textbf{0.996} & \textbf{0.935} \\ \hline
\end{tabular}
\end{table}

\begin{figure*}[t]
\centering
\figph[0.95\textwidth]{confusion matrices of the three baselines}{fig-confusion}
\caption{Confusion matrices of the three baselines on the test set (threshold selected on validation).}
\label{fig:cm}
\end{figure*}

Gradient Boosting detects 257 of the 290 positive samples with 51 false alarms
(TN = 4,989, FP = 51, FN = 33, TP = 257), followed by Random Forest (247 detected, 57 false
alarms) and Logistic Regression (204 detected, 101 false alarms). The accuracy of all models
(0.965 to 0.984) is above the 0.945 of the trivial classifier but hides the differences in the
minority class, which confirms that accuracy is not informative here.

The decision threshold matters. At the default 0.5, Logistic Regression reaches a recall of
0.962 but a precision of only 0.361 (F1 = 0.525), i.e., many false alarms; the threshold
selected on validation (0.932) trades part of that recall for a much better F1 (0.686).

\subsection{Ablation: Feature Engineering and Augmentation}
Table~\ref{tab:ablation} changes one factor at a time with respect to the main setting
(engineered features, no augmentation).

\begin{table}[t]
\caption{ABLATION: F1 / PR-AUC ON THE TEST SET}
\label{tab:ablation}
\centering
\begin{tabular}{|l|c|c|c|}
\hline
\textbf{Model} & \textbf{Raw feat.} & \textbf{Engineered} & \textbf{Eng. + aug.} \\ \hline
Logistic Regression & 0.230 / 0.098 & 0.686 / 0.791 & 0.685 / 0.792 \\ \hline
Random Forest & 0.839 / 0.944 & 0.832 / 0.923 & 0.826 / 0.933 \\ \hline
Gradient Boosting & 0.873 / 0.952 & 0.860 / 0.935 & 0.848 / 0.942 \\ \hline
\end{tabular}
\end{table}

\emph{Feature engineering.} For Logistic Regression the engineered features raise the ROC-AUC
from 0.542 (almost random) to 0.981 and the PR-AUC from 0.098 to 0.791: the raw decision
boundary is not linear, but distances to the failure limits make it nearly linear.
For the tree ensembles the raw features are already enough, and the differences between
raw and engineered inputs (up to 0.013 in F1) are small relative to what a single split of
60 test episodes can resolve, so we do not claim that either is better for these models.

\emph{Augmentation.} Noise augmentation did not improve any model (F1 differences between
$-0.012$ and $-0.001$). With a clean simulator, the added noise does not seem to provide
useful variety.

\subsection{Problem Difficulty}
Tree ensembles reach a ROC-AUC close to 0.995, so the classes are almost separable in a
ranking sense, yet the F1 stays around 0.83--0.87: the remaining errors concentrate on the
precision/recall trade-off of the rare class. A plausible reason is that states a few steps
before the label window are almost identical to states inside it (the label boundary is
arbitrary with respect to the dynamics), but we did not verify this and leave it as a
hypothesis. The strong results also leave limited headroom for the deep learning models of
the next workshop, which will have to be compared against a high baseline.

\subsection{Limitations}
All results come from one episode-level split with a single seed and 60 test episodes, so
small differences (e.g., Random Forest versus Gradient Boosting) are not conclusive.
Cross-validation grouped by episode and several seeds would give variance estimates.
Hyperparameters were not tuned beyond sensible defaults, and the dataset covers only the
behaviour of the mixed policy (noisy expert plus random actions) that generated it.

%% ======================================================================
\section{Conclusions}
We characterised a CartPole failure-prediction dataset, identified a 17.2:1 imbalance and
the need to split by episode, and built a reproducible pipeline with scaling and
augmentation restricted to the training data. Gradient Boosting was the strongest baseline
(F1 = 0.860, recall = 0.886, PR-AUC = 0.935), Random Forest was close, and Logistic
Regression required engineered distance-to-limit features to be useful. Noise augmentation
brought no gain. These numbers are the reference for the next workshops. Future work
includes grouped cross-validation, a multi-seed evaluation, and sequence models that
exploit the temporal structure of the episodes.

\section*{Reproducibility}
Code, split file and instructions are available at \url{https://github.com/Tonyaa11/Cartpole}. All
reported values are produced by \texttt{src/baselines.py} (seed 42) in Google Colab
(scikit-learn 1.6.1, pandas 2.2.3, numpy 2.1.3); other library versions can change the
Random Forest results slightly.

\begin{thebibliography}{00}
\bibitem{barto1983} A. G. Barto, R. S. Sutton, and C. W. Anderson, ``Neuronlike adaptive elements that can solve difficult learning control problems,'' \emph{IEEE Trans. Syst., Man, Cybern.}, vol. SMC-13, no. 5, pp. 834--846, 1983.
\bibitem{pedregosa2011} F. Pedregosa et al., ``Scikit-learn: Machine learning in Python,'' \emph{J. Mach. Learn. Res.}, vol. 12, pp. 2825--2830, 2011.
\bibitem{breiman2001} L. Breiman, ``Random forests,'' \emph{Mach. Learn.}, vol. 45, no. 1, pp. 5--32, 2001.
\bibitem{friedman2001} J. H. Friedman, ``Greedy function approximation: A gradient boosting machine,'' \emph{Ann. Statist.}, vol. 29, no. 5, pp. 1189--1232, 2001.
\bibitem{saito2015} T. Saito and M. Rehmsmeier, ``The precision-recall plot is more informative than the ROC plot when evaluating binary classifiers on imbalanced datasets,'' \emph{PLoS ONE}, vol. 10, no. 3, e0118432, 2015.
\end{thebibliography}

\end{document}
